\documentclass[11pt]{article}
\usepackage[a4paper,margin=27mm]{geometry}
\usepackage{fontspec}
\usepackage{amsmath,amsthm,mathtools}
\usepackage{unicode-math}
\setmainfont[
  UprightFont=texgyrepagella-regular.otf,
  BoldFont=texgyrepagella-bold.otf,
  ItalicFont=texgyrepagella-italic.otf,
  BoldItalicFont=texgyrepagella-bolditalic.otf
]{texgyrepagella-regular.otf}
\setsansfont[
  UprightFont=texgyreheros-regular.otf,
  BoldFont=texgyreheros-bold.otf
]{texgyreheros-regular.otf}
\setmonofont{texgyrecursor-regular.otf}
\setmathfont{texgyrepagella-math.otf}
\newfontfamily\zhfont[
  Path=assets/fonts/,
  UprightFont=NotoSerifCJKtc-Regular.otf
]{NotoSerifCJKtc-Regular.otf}
\newfontfamily\zhfontbold[
  Path=assets/fonts/,
  UprightFont=NotoSerifCJKtc-Regular.otf,
  FakeBold=2
]{NotoSerifCJKtc-Regular.otf}
\usepackage{booktabs,array,enumitem,microtype,xcolor}
\usepackage[authoryear,round]{natbib}
\usepackage[hidelinks]{hyperref}
\usepackage{fancyhdr}
\hypersetup{
  pdftitle={Pinching a free normal cover: variable critical exponent in reduced quasiconformal Teichmuller space},
  pdfauthor={Anonymous candidate contribution},
  pdfsubject={Unrefereed all-Fuchsian strengthening candidate for AIM Problem 4.3},
  pdfkeywords={Fuchsian group, critical exponent, quasiconformal Teichmuller space, normal cover, bottom of spectrum, collar lemma}
}

\definecolor{statusred}{RGB}{137,34,44}
\pagestyle{fancy}
\fancyhf{}
\fancyhead[L]{\small\itshape Unrefereed strengthening candidate}
\fancyhead[R]{\small\itshape AIM Problem 4.3}
\fancyfoot[C]{\thepage}
\setlength{\headheight}{14pt}
\setlength{\emergencystretch}{2em}

\newtheorem{theorem}{Theorem}
\newtheorem{lemma}[theorem]{Lemma}
\newtheorem{proposition}[theorem]{Proposition}
\newtheorem{corollary}[theorem]{Corollary}
\theoremstyle{definition}
\newtheorem{definition}[theorem]{Definition}
\newtheorem{remark}[theorem]{Remark}

\newcommand{\Hh}{\mathbb H^2}
\newcommand{\Fg}{F_g}
\newcommand{\Tqcred}{T_{\mathrm{qc}}^{\mathrm{red}}}
\newcommand{\lamzero}{\lambda_0}

\title{\textbf{Pinching a free normal cover:\\
variable critical exponent in reduced quasiconformal Teichm\"uller space}}
\author{Anonymous candidate contribution}
\date{1 September 2026}

\begin{document}
\maketitle

\begin{abstract}
We give an explicit all-Fuchsian strengthening candidate related to
M.~Kapovich's AIM Problem~4.3.
For every closed oriented surface $S_g$ of genus $g\ge2$, fix the kernel $K$
of the epimorphism $\pi_1(S_g)\to F_g$ that kills a disjoint cut system.
For every marked closed hyperbolic metric $m$, the image $\Gamma_m$ of $K$
under the Fuchsian holonomy is an infinitely generated Fuchsian group of the
first kind.  Marked quasiconformal maps of the compact base lift to
quasiconformal conjugacies of the groups $\Gamma_m$, so this gives a family in
one reduced, surface-based quasiconformal Teichm\"uller space.  The quotient
$F_g$ is nonamenable;
the normal-subgroup theorem of Dougall and Sharp therefore gives
$\delta(\Gamma_m)<1$ for every marked compact-base metric.  If the cut curves are
simultaneously pinched to common length $\ell$, a compactly supported function
on one lifted cut-surface cell has Rayleigh quotient
$O_g(\ell)$.  The Elstrodt--Patterson--Sullivan formula then yields
\[
 \frac12<\delta(\Gamma_{m_\ell})<1,
 \qquad
 1-\delta(\Gamma_{m_\ell})=O_g(\ell),
 \qquad
 \delta(\Gamma_{m_\ell})\longrightarrow1.
\]
Thus the critical exponent assumes infinitely many distinct values among the
all-Fuchsian points of $\Tqcred(X_0)$.  Under the reduced surface-based reading
of the AIM question this is a direct solution candidate; under the broader
group-deformation reading it is an all-Fuchsian strengthening of an affirmative
precedent of Astala and Zinsmeister.
\end{abstract}

\begingroup
\XeTeXlinebreaklocale ``zh''
\XeTeXlinebreakskip=0pt plus 1pt
\zhfont
\par\bigskip
{\zhfontbold\large 中文摘要（繁體）\par}
\smallskip
本文提出 M.~Kapovich 之 AIM 問題 4.3 的一個明確解答候選。對每個虧格
$g\ge2$ 的閉可定向曲面 $S_g$，固定滿同態
$\pi_1(S_g)\to F_g$ 的核 $K$；此滿同態消去一組互不相交的割曲線。
對任一有標記的閉雙曲度量 $m$，$K$ 在 Fuchsian 單值群中的像
$\Gamma_m$ 是第一類且無限生成的 Fuchsian 群。緊緻底曲面之間的有標記
擬共形映射可提升為 $\Gamma_m$ 之間的擬共形共軛，故所有這些群位於同一
約化擬共形 Teichm\"uller 空間。商群 $F_g$ 非可均，因此 Dougall--Sharp 定理給出每一
有限階段的 $\delta(\Gamma_m)<1$。若同時把割曲線壓縮至共同長度 $\ell$，
則一個提升之割開胞腔上的緊支撐函數具有
$O_g(\ell)$ 的 Rayleigh 商。Elstrodt--Patterson--Sullivan
公式遂推出 $1/2<\delta(\Gamma_{m_\ell})<1$，且
$\delta(\Gamma_{m_\ell})\to1$。因此臨界指數在同一擬共形變形空間中取得
無限多個不同值。
\endgroup

\smallskip
\noindent\emph{Translation note.} The Traditional Chinese abstract is a
producer-generated rendering and was not independently assessed by a Chinese
language reviewer; the English abstract controls if the two differ.

\medskip
\noindent\colorbox{statusred!8}{\parbox{0.95\linewidth}{
\textcolor{statusred}{\textbf{Status.}} This is an anonymous, AI-assisted,
unrefereed all-Fuchsian strengthening candidate.  It includes a complete written proof and
producer-side formula checks, but no claim of historical priority, external
specialist validation, independent reconstruction, formal proof verification,
or peer review.  Candidate publication and DOI assignment do not alter those
assurance limits.}}

\medskip
\noindent\textbf{2020 Mathematics Subject Classification.}
Primary 30F35; Secondary 30F60, 58J50.\quad
\textbf{Keywords.} Fuchsian group, critical exponent, quasiconformal
Teichm\"uller space, normal cover, bottom of spectrum, collar degeneration.

\tableofcontents

\section{Introduction}

Let $\Gamma_0$ be an infinitely generated Fuchsian group of the first kind.
AIM Problem~4.3, attributed to M.~Kapovich, states:
\begin{quote}
``Construct an infinitely generated Fuchsian group $\Gamma_0$ of the first kind
such that the critical exponents of some elements in $T_{\rm qc}(\Gamma_0)$ are
different.''
\end{quote}
The problem arose at the 2012 AIM workshop, appears in the accessible AIM
problem document dated 2014, and was published in handbook form in 2016
\citep{AIM2014,Su2016}.  This question is meaningful because first kind means that the
limit set is the entire boundary circle, whereas for an infinitely generated
group the critical exponent can nevertheless be strictly smaller than one.
The critical exponent measures exponential orbit growth: roughly, how rapidly
the number of group translates within distance $R$ grows with $R$.  A first-kind
group has the full boundary circle as limit set, but for an infinitely generated
group this topological fact does not force exponent one.  Pinching the cut
curves creates long collars, so a function concentrated on one lifted cell has
very low spectral energy; the spectral formula then drives the exponent toward
one.  Examples with exponent below one occur elsewhere
\citep{Bishop2003,Lehnert2017}.

There is an important scope boundary.  Under a broad group-deformation reading,
the quasi-Fuchsian family of \citet{AstalaZinsmeister1995} appears already to
give an affirmative existential example.  We therefore make no first-solution
or priority claim.  The contribution here is a fixed-kernel construction that
stays entirely in the all-Fuchsian locus, has every nondegenerate exponent below
one, and approaches one with an explicit one-sided rate.  Under the narrower
reduced surface-based, all-Fuchsian reading of the AIM question, it is a direct
solution candidate.

The construction here uses a fixed normal cover of a closed surface.  Its deck
group is a nonabelian free group.  Two competing consequences of that choice
drive the proof.  At every fixed metric, nonamenability makes the normal
subgroup's critical exponent strictly smaller than the ambient lattice
exponent.  On the other hand, when a complete cut system is pinched, one
fundamental cell becomes separated from its neighbours by increasingly long
collars.  A function supported on that cell has arbitrarily small Dirichlet
energy.  The spectral formula converts this geometric degeneration into
critical exponents tending to one.

We use the nonnegative Laplacian.  For a complete surface $X$, its spectral
bottom is
\[
 \lamzero(X)=\inf_{0\ne f\in C_c^\infty(X)}
 \frac{\int_X |\nabla f|^2\,dA}{\int_X |f|^2\,dA}.
\]
For a nonelementary Fuchsian group $\Gamma$, the critical exponent is
\[
 \delta(\Gamma)=\inf\left\{s>0:
 \sum_{\gamma\in\Gamma}e^{-s d(o,\gamma o)}<\infty\right\},
\]
independent of $o\in\Hh$.

Fix a base metric $m_0$ and put $X_0=X_{m_0}$.  We use the reduced,
surface-based convention of \citet{AlessandriniEtAl2012}: a point of
$\Tqcred(X_0)$ is an equivalence class of a marked conformal surface $(Y,f)$,
where $f:X_0\to Y$ is quasiconformal, and $(Y_1,f_1)$ and $(Y_2,f_2)$ are
equivalent when a conformal map $c:Y_1\to Y_2$ is homotopic to
$f_2\circ f_1^{-1}$.  In the reduced convention boundary components are
preserved setwise rather than pointwise; the surfaces here have no boundary,
so that distinction is moot.  The theorem concerns the Fuchsian uniformising
groups of the marked surfaces constructed below.  It does not quantify over
arbitrary quasi-Fuchsian sphere conjugates.

For comparison with the original notation, write $T_{\rm qc}(\Gamma_0)$ for
the reduced group quasiconformal deformation space: a quasiconformal sphere
homeomorphism $W$ conjugating $\Gamma_0$ to a Kleinian group determines a
point, and marked representatives that differ by M\"obius conjugacy determine
the same point.  This is the only group-space structure used below.  M\"obius
conjugacy preserves the critical exponent, since it acts isometrically on
hyperbolic three-space.  Lemma~\ref{lem:qc} constructs such a sphere
representative for every marked compact-base metric and shows that the entire
family lies in the all-Fuchsian locus of $T_{\rm qc}(\Gamma_0)$.  Thus the
construction produces points for the AIM notation without asserting a global
identification of every group deformation with the surface-based space.

\section{The fixed free normal cover}

Fix $g\ge2$.  Choose a geometric symplectic basis
$(a_1,b_1,\ldots,a_g,b_g)$ for $S_g$ so that the simple curves
$b_1,\ldots,b_g$ are pairwise disjoint and their complement is connected.
Thus
\[
 \pi_1(S_g)=\left\langle a_1,b_1,\ldots,a_g,b_g\ \middle|\
 \prod_{i=1}^g[a_i,b_i]=1\right\rangle .
\]
Let $F_g=\langle x_1,\ldots,x_g\rangle$ and define
\begin{equation}\label{eq:q}
 q:\pi_1(S_g)\longrightarrow F_g,
 \qquad q(a_i)=x_i,\qquad q(b_i)=1.
\end{equation}
The surface relator maps to the identity, and the images of the $a_i$ generate
$F_g$; hence $q$ is an epimorphism.  Put $K=\ker q$.
Thus there is an exact sequence
\[
 1\longrightarrow K\longrightarrow\pi_1(S_g)
 \mathrel{\mathop{\longrightarrow}^{q}}F_g\longrightarrow1.
\]

For a marked closed hyperbolic metric $m$ on $S_g$, choose its faithful
Fuchsian holonomy
\[
 \rho_m:\pi_1(S_g)\longrightarrow \mathrm{PSL}(2,\mathbb R),
\]
write $L_m=\rho_m(\pi_1(S_g))$, and set
\[
 \Gamma_m=\rho_m(K),\qquad X_m=\Hh/\Gamma_m.
\]
Then $X_m\to S_m=\Hh/L_m$ is the regular cover with deck group $F_g$.
Equivalently,
\[
 1\longrightarrow\Gamma_m\longrightarrow L_m
 \longrightarrow F_g\longrightarrow1.
\]

\begin{lemma}\label{lem:first-kind}
For every $m$, the group $\Gamma_m$ is nonelementary, infinitely generated,
and of the first kind.
\end{lemma}

\begin{proof}
The subgroup $K$ is nontrivial because it contains each $b_i$.  Hence
$\Gamma_m$ contains a hyperbolic element.  Its limit set is nonempty and,
because $\Gamma_m\triangleleft L_m$, is invariant under $L_m$.  The limit set
of the cocompact lattice $L_m$ is the full circle and is the minimal nonempty
closed $L_m$-invariant subset of the boundary.  It follows that
\[
 \Lambda(\Gamma_m)=\Lambda(L_m)=\partial\Hh.
\]
In particular $\Gamma_m$ is nonelementary and of the first kind.

Suppose that $\Gamma_m$ were finitely generated.  A finitely generated
Fuchsian group without parabolics is convex cocompact
\citep[see, for example,][]{Beardon1983}.  There are no parabolics here because
$\Gamma_m$ is a subgroup of the cocompact torsion-free group $L_m$.  Since its
limit set is the full circle, its convex hull is all of $\Hh$; convex
cocompactness would therefore make $X_m=\Hh/\Gamma_m$ compact.  A compact
cover of the compact surface $S_m$ has finite degree.  This contradicts
$L_m/\Gamma_m\cong F_g$, which is infinite.  Thus $\Gamma_m$ is infinitely
generated.
\end{proof}

\begin{lemma}\label{lem:qc}
Fix a marked metric $m_0$.  For every marked closed hyperbolic metric $m$,
the lifted marked surface $X_m$ represents a point of $\Tqcred(X_0)$, and its
Fuchsian uniformising group is $\Gamma_m$.  Consequently the family
$\{\Gamma_m:m\in\operatorname{Teich}(S_g)\}$ is a family of Fuchsian points in
one reduced quasiconformal Teichm\"uller space.
\end{lemma}

\begin{proof}
There is an orientation- and marking-preserving quasiconformal map
$h:S_{m_0}\to S_m$.  Choose the lift
$\widetilde h:\Hh\to\Hh$ that is equivariant for the marked holonomies.  Then
\[
 \widetilde h\,\rho_{m_0}(\gamma)\,\widetilde h^{-1}
   =\rho_m(\gamma)\qquad(\gamma\in\pi_1(S_g)).
\]
Restricting this identity to the fixed subgroup $K$ gives
$\widetilde h\Gamma_{m_0}\widetilde h^{-1}=\Gamma_m$.  The lift is
quasiconformal and descends to the marked map $X_0\to X_m$, so it gives the
asserted point of $\Tqcred(X_0)$.  The markings identify the same abstract
kernel $K=\ker q$ at both ends.  Its boundary extension is quasisymmetric,
and reflection across the boundary circle supplies a quasiconformal sphere
extension $W_m$ satisfying
$W_m\Gamma_{m_0}W_m^{-1}=\Gamma_m$.  Hence it also gives a point of
$T_{\rm qc}(\Gamma_{m_0})$ in the original AIM notation.  If two surface
representatives are equivalent, the conformal equivalence lifts to a M\"obius
conjugacy of their uniformising groups; consequently their critical exponents
agree.  That observation
does not enlarge the theorem to arbitrary quasi-Fuchsian conjugates.  No
injectivity claim for the induced map of Teichm\"uller spaces is needed.
\end{proof}

\section{Strict comparison with the ambient lattice}

The ambient group $L_m$ is cocompact, hence convex cocompact, and
$\delta(L_m)=1$.  The quotient
\[
 L_m/\Gamma_m\cong F_g
\]
is nonamenable because $g\ge2$.  The normal-subgroup theorem of
\citet[Theorem~1.1]{DougallSharp2016} says that for a normal subgroup of a
convex-cocompact group acting on a pinched Hadamard manifold, equality of
critical exponents holds exactly when the quotient is amenable.  Here
$\Hh$ is complete, simply connected, and has curvature $-1$; $L_m$ is
nonelementary and cocompact, hence convex cocompact; $\Gamma_m\triangleleft
L_m$; and $L_m/\Gamma_m\cong F_g$ is nonamenable.  Every hypothesis is
therefore discharged, and
\begin{equation}\label{eq:strict}
 \delta(\Gamma_m)<\delta(L_m)=1
 \qquad\text{for every marked metric }m.
\end{equation}

This is the only place where nonamenability is used.  Equation
\eqref{eq:strict} also chooses the correct branch when the spectral formula is
applied after degeneration.

\section{Pinching the cut system}

Extend $b_1,\ldots,b_g$ to a pants decomposition.  Fenchel--Nielsen
coordinates give a path of marked closed hyperbolic metrics $m_\ell$,
$0<\ell<\ell_*$, for which
\[
 \ell_{m_\ell}(b_i)=\ell\qquad(1\le i\le g).
\]
The remaining length and twist coordinates may be held fixed
\citep{Buser1992}.  Let
$X_\ell=X_{m_\ell}$ and $\Gamma_\ell=\Gamma_{m_\ell}$.

Cutting $S_g$ along the $b_i$ produces a sphere $P$ with $2g$ boundary
components.  The cover defined by \eqref{eq:q} can be constructed from copies
of $P$ indexed by $F_g$: after orienting the two sides of each $b_i$, the
corresponding sides of the copies labelled $h$ and $hx_i$ are glued.  Indeed,
the curve $a_i$ is dual to $b_i$: after the cut it becomes an arc joining the
two boundary copies of $b_i$, crosses no other cut curve, and its monodromy is
$q(a_i)=x_i$.  In
particular, $X_\ell$ contains an embedded compact lifted cell $P_\ell$.

The collar lemma supplies pairwise disjoint open collars around the $b_i$ with
half-width
\begin{equation}\label{eq:w}
 w=w(\ell)=\operatorname{arcsinh}\!\left(
 \frac{1}{\sinh(\ell/2)}\right)
 \longrightarrow\infty
 \qquad(\ell\downarrow0)
\end{equation}
and metric
\begin{equation}\label{eq:collar-metric}
 ds^2=dr^2+\ell^2\cosh^2r\,d\theta^2,
 \qquad -w<r<w,\quad \theta\in\mathbb R/\mathbb Z
\end{equation}
\citep{Keen1974,Buser1992}.  Each cut cell contains $2g$ half-collars,
parametrized by $0\le r\le w$, with the cut geodesic at $r=0$.

For sufficiently small $\ell$, $w(\ell)>1$.  Define
$u_\ell:X_\ell\to[0,1]$ as follows.  On each of the $2g$ selected-cell
half-collars put
\[
 u_\ell(r,\theta)=
 \begin{cases}
  r,&0\le r\le1,\\
  1,&1\le r<w(\ell).
 \end{cases}
\]
Put $u_\ell=1$ on the remainder of that cell and $u_\ell=0$ outside it.  The
traces agree across every seam because $u_\ell=0$ at $r=0$.  Thus $u_\ell$ is
a compactly supported Lipschitz
function.  It belongs to $H^1_c(X_\ell)$ and can be approximated in $H^1$ by
functions in $C_c^\infty(X_\ell)$, so it is admissible in the Rayleigh
principle.

\begin{proposition}\label{prop:lambda}
There are constants $C_g>0$ and $\ell_g>0$ such that
\[
 0\le \lamzero(X_\ell)\le
 C_g\ell
 \qquad(0<\ell<\ell_g).
\]
In particular $\lamzero(X_\ell)\to0$.
\end{proposition}

\begin{proof}
The gradient is supported in $0<r<1$, where
$|\nabla u_\ell|^2=1$.  The area element from
\eqref{eq:collar-metric} is $\ell\cosh r\,dr\,d\theta$.  One half-collar
therefore contributes
\[
 \int_0^1\!\int_0^1 \ell\cosh r\,dr\,d\theta
 =\ell\sinh1.
\]
Summing over $2g$ half-collars gives the exact numerator
\begin{equation}\label{eq:energy}
 E_\ell=\int_{X_\ell}|\nabla u_\ell|^2\,dA
 =2g\ell\sinh1.
\end{equation}

By Gauss--Bonnet, $\operatorname{area}(S_{m_\ell})=4\pi(g-1)$.  Each of the
$2g$ width-one transition strips has area $\ell\sinh1$.  Hence the plateau on
which $u_\ell=1$ has area at least
\begin{equation}\label{eq:core-area}
 D_\ell=4\pi(g-1)-2g\ell\sinh1.
\end{equation}
For all sufficiently small $\ell$, specifically whenever
$2g\ell\sinh1\le2\pi(g-1)$, we have
$D_\ell\ge2\pi(g-1)>0$, and therefore
\[
 \int_{X_\ell}u_\ell^2\,dA\ge D_\ell\ge2\pi(g-1).
\]
The Rayleigh principle and \eqref{eq:energy} give
\[
 \lamzero(X_\ell)
 \le\frac{E_\ell}{\int u_\ell^2}
 \le\frac{g\sinh1}{\pi(g-1)}\,\ell.
\]
Thus one may take $C_g=g\sinh1/[\pi(g-1)]$ after decreasing $\ell_g$ if
necessary.
\end{proof}

\section{Critical exponents and the all-Fuchsian conclusion}

For any nonelementary Fuchsian group $\Gamma$, the
Elstrodt--Patterson--Sullivan formula, in the nonnegative-Laplacian convention,
is
\begin{equation}\label{eq:eps}
 \lamzero(\Hh/\Gamma)=
 \begin{cases}
  \tfrac14, & \delta(\Gamma)\le\tfrac12,\\[2mm]
  \delta(\Gamma)(1-\delta(\Gamma)), & \delta(\Gamma)\ge\tfrac12.
 \end{cases}
\end{equation}
See \citet[Theorem~2.17, p.~333]{Sullivan1987}.  Its assumptions are met:
$\Gamma_\ell$ is torsion-free, discrete, and nonelementary; $X_\ell$ is a
complete curvature-$-1$ surface; and our $\lamzero$ uses the nonnegative
Laplacian.  We can now state the result.

\begin{theorem}[All-Fuchsian strengthening candidate for AIM Problem 4.3]\label{thm:main}
For every $g\ge2$, the group $\Gamma_0=\Gamma_{m_0}$ constructed above is an
infinitely generated Fuchsian group of the first kind for which the critical
exponent takes infinitely many distinct values among the Fuchsian points of
$\Tqcred(X_0)$.  More
precisely, along the simultaneous pinching family, for all sufficiently small
$\ell$,
\[
 \frac12<\delta(\Gamma_\ell)<1,
 \qquad
 0<1-\delta(\Gamma_\ell)
 \le 2C_g\ell,
\]
and $\delta(\Gamma_\ell)\to1$ as $\ell\downarrow0$.
\end{theorem}

\begin{proof}
Lemmas~\ref{lem:first-kind} and \ref{lem:qc} give the required group and place
the whole family in $\Tqcred(X_0)$.  By
Proposition~\ref{prop:lambda}, $\lamzero(X_\ell)<1/4$ for all sufficiently
small $\ell$.  Formula \eqref{eq:eps} then rules out
$\delta(\Gamma_\ell)\le1/2$, so
\[
 \lamzero(X_\ell)=
 \delta(\Gamma_\ell)(1-\delta(\Gamma_\ell)).
\]
Equation \eqref{eq:strict} supplies $\delta(\Gamma_\ell)<1$.  Consequently
\[
 0<1-\delta(\Gamma_\ell)
 =\frac{\lamzero(X_\ell)}{\delta(\Gamma_\ell)}
 <2\lamzero(X_\ell)
 \le2C_g\ell.
\]
This proves $\delta(\Gamma_\ell)\to1$ from below.  No exponent belonging to a
marked compact-base metric is one, so the values along any sequence
$\ell_n\downarrow0$ cannot form a finite
set: a finite set contained in $(-\infty,1)$ has maximum strictly smaller than
one.  Thus infinitely many values are distinct, and in particular the
critical exponent is nonconstant among the Fuchsian points of
$\Tqcred(X_0)$.
\end{proof}

\begin{remark}
The proof does not show that $\delta(\Gamma_\ell)$ is monotone in $\ell$.
It also does not identify an exact formula for the exponent.  The explicit
rate is one-sided and is inherited from a single fixed-width test function.
The symbol $\ell=0$ denotes a boundary degeneration, not a metric or a point
of the asserted family.
\end{remark}

\section{Relation to previous work}

\citet{AstalaZinsmeister1995} used a normal cover of a compact genus-three
surface with abelian deck group $\mathbb Z^3$ and holomorphic quasi-Fuchsian
deformations to exhibit variation of the Poincar\'e exponent, including values
above one.  Under the broad group-deformation interpretation of
$T_{\rm qc}(\Gamma_0)$, this appears already to answer the existential AIM
question affirmatively.  The present theorem is therefore positioned as an
all-Fuchsian strengthening and alternative construction, and as a direct
solution candidate only under the reduced surface-based, all-Fuchsian reading.
It is not the same family: here the deck group is the
nonamenable free group $F_g$, every compact-base exponent is below one, and
pinching makes the exponents approach one from below.  Bishop's
$\delta$-stability results and the later Ruelle-property work concern the
relation between exponent, limit-set dimension, and quasi-Fuchsian
deformation \citep{Bishop2003,HuoZinsmeister2022}.  Regular covers and the
sharp lower threshold $1/2$ for their exponents are studied by
\citet{BonfertTaylorEtAl2012}.  These results delimit the contribution: the
claim here is the fixed-kernel, all-Fuchsian pinching construction in reduced
surface-based quasiconformal Teichm\"uller space, not a general deformation or
stability theorem.

\section{Dependency and limitation statement}

The family construction, quasiconformal lift, collar function, energy
integral, area estimate, and final algebra are given explicitly.  The proof
imports two decisive structural results: the Dougall--Sharp normal-subgroup
amenability criterion and the Elstrodt--Patterson--Sullivan spectral formula.
Dougall--Sharp Theorem~1.1 is applied with $X=\Hh$, $\Gamma=L_m$, and
$\Gamma_0=\Gamma_m$; Sullivan Theorem~2.17 is applied in dimension $d=1$ to
$\Hh/\Gamma_\ell$.  The hypothesis maps are also recorded in
\texttt{PROOF\_OBLIGATIONS.md}.  The argument uses standard finite-type
Teichm\"uller and collar theorems.  The
accompanying program recomputes only the explicit collar formulas for sample
parameters; it does not verify these external theorems or replace the
universal written argument.

A bounded search through the AIM source, arXiv, publisher and author-hosted
primary papers located the adjacent work just described but did not locate
this exact fixed free-normal-cover, all-Fuchsian pinching statement.  That
negative search evidence is not a priority claim.  Specialist review and
exhaustive MathSciNet/zbMATH searching remain open assurance steps.

\section*{Data and materials availability}

There are no empirical data.  The local research package contains the LaTeX
and Markdown manuscripts, bibliography, source and claim ledgers,
standard-library formula checker, mutation tests, build scripts, PDF, and
cryptographic manifest.  The public repository is
\url{https://github.com/ipitchford/pinching-free-normal-cover-critical-exponent},
and the version DOI is \url{https://doi.org/10.5281/zenodo.22229561}.
The PDF is visually checked but untagged; the Markdown
manuscript is the accessible-text fallback, with the remaining limitation that
its mathematics is linearized plain text rather than fully semantic math.
Their existence does not increase the mathematical assurance level.

\section*{Ethics statement}

This theoretical mathematics work uses no human participants, personal data,
animals, clinical material, or field interventions.  No ethics approval was
required.

\section*{Author contributions (CRediT)}

Anonymous candidate contribution: Conceptualization, Formal analysis,
Methodology, Software, Validation, Writing---original draft,
and Writing---review \& editing.  Operational provenance and AI assistance are
reported separately below; no AI system is listed as an author.

\section*{Funding}

No external funding was declared for this candidate.

\section*{Competing interests}

No competing interests were declared.

\section*{AI-use disclosure}

OpenAI Codex was used to search and inspect sources, develop and stress-test
the construction, draft and typeset the manuscript, write corroborative
checks, and assemble the local package.  The workflow is producer-side and
does not constitute independent mathematical validation.  A human publisher
or reviewer should verify the proof, citations, attribution, and novelty
boundary before any public release.

\bibliographystyle{plainnat}
\bibliography{references}

\end{document}
